\documentclass[11pt]{amsart}
\usepackage[margin=1in]{geometry}
\usepackage{amssymb,mathtools,booktabs}
\usepackage{xurl}
\usepackage[colorlinks=true,linkcolor=blue,citecolor=blue,urlcolor=blue]{hyperref}
\newtheorem{theorem}{Theorem}[section]
\newtheorem{proposition}[theorem]{Proposition}
\newtheorem{lemma}[theorem]{Lemma}
\theoremstyle{remark}
\newtheorem{remark}[theorem]{Remark}
\newcommand{\F}{\mathbb F}

\title[{Jacobsthal's function in $\F_q[x]$: the first two values}]{A remark on Jacobsthal's function in $\F_q[x]$: \\ the first two values}
\author{[AUTHOR NAME]}
\date{October 9, 2026}

\begin{document}
\begin{abstract}
OpenAI's preprint \emph{A quadratic bound for Jacobsthal's function} (family 021 of the OpenAI
mathematics repository) proves $h(k)\ll k^2/(\log\log 3k)^2$ for the integer Jacobsthal function.
We record the elementary first cases of the analogous quantity in the polynomial ring $\F_q[x]$.
Let $K_q(m)$ be the least number of distinct monic irreducibles of a polynomial $M$ for which some
``interval'' $f+\{g:\deg g<m\}$ contains no polynomial coprime to $M$. We show $K_q(1)=q$ and
$K_q(2)=2q-1$ for every prime power $q$, characterise the extremal configurations for $q\ge3$, and
list exact values for a few further small cases computed in two independent ways. The arguments are
short counting arguments; nothing from the integer proof is used.
\end{abstract}
\maketitle

\noindent\textbf{Attribution and scope.} The integer problem and its solution are OpenAI's work
(family 021 of \url{https://github.com/openai/math}, commit
\nolinkurl{fd4aeeb2ee4fc729c18d98444fed42fd0529eeeb}) \cite{oai}. This note is only a function-field
analogue in the smallest cases. It says nothing about the integer problem and nothing about the
asymptotics of $K_q(m)$ in $m$. Developed with AI assistance (Anthropic Claude, Opus~5.5) under
[AUTHOR NAME]'s direction; the independent adversarial review was run by a separate AI session and
is not human peer review.

\section{Definitions}
Let $q$ be a prime power and $V_m=\{g\in\F_q[x]:\deg g<m\}$ (so $|V_m|=q^m$; $\deg 0=-\infty$).
For $M\ne0$ let $j(M)$ be the least $m\ge0$ such that every coset $f+V_m$, $f\in\F_q[x]$, contains
an element coprime to $M$, and put $h_q(k)=\max\{j(M):\omega(M)\le k\}$, where $\omega$ counts
distinct monic irreducible factors. Dually,
\[
 K_q(m)=\min\{\omega(M): \exists f\ \forall g\in V_m:\ \gcd(f+g,M)\ne1\},
\]
so that $h_q(k)\le m$ if and only if $K_q(m)>k$. (The zero polynomial is not coprime to $M$, and
nonzero constants are coprime to every $M$; both conventions are as for integers.)

\begin{lemma}[covering model]\label{lem:model}
Fix $m\ge1$. $K_q(m)$ equals the minimum, over sets $\mathcal P$ of monic irreducibles of degree $<m$
with a residue $r_P$ chosen for each $P\in\mathcal P$, of
$|\mathcal P|+\#\{g\in V_m: g\not\equiv r_P \pmod P\ \forall P\in\mathcal P\}$.
\end{lemma}
\begin{proof}
If $\deg P\ge m$ then $P$ divides at most one element of $f+V_m$, since the difference of two
distinct elements is nonzero of degree $<m$. If $\deg P<m$ then $P\mid f+g$ iff $g\equiv-f\pmod P$.
Hence any admissible $M$ for $f$ gives a configuration of the stated cost, with $r_P=-f\bmod P$.
Conversely, given a configuration and the set $U$ of uncovered points, choose distinct monic
irreducibles $Q_u$ ($u\in U$) of degree $\ge m$ not in $\mathcal P$, and by the Chinese remainder
theorem choose $f$ with $f\equiv-r_P\pmod P$ for $P\in\mathcal P$ and $f\equiv-u\pmod{Q_u}$ for
$u\in U$. Then $M=\prod_{\mathcal P}P\prod_U Q_u$ works and $\omega(M)$ equals the cost.
\end{proof}

\section{The first two values}
\begin{theorem}\label{thm:main}
For every prime power $q$: $K_q(1)=q$ and $K_q(2)=2q-1$. Equivalently $h_q(k)=1$ for $1\le k\le q-1$,
$h_q(k)=2$ for $q\le k\le 2q-2$, and $h_q(2q-1)\ge3$. For $q\ge3$, a configuration attains
$K_q(2)$ if and only if it uses all $q$ linear primes with residues $r_{x-a}=g^*(a)$
$(a\in\F_q)$ for some $g^*\in V_2$ (equivalently, after replacing $f$ by an element of the same
coset, $x^q-x$ divides $f$), and in addition exactly $q-1$ primes of degree $\ge2$, each dividing a
different one of the $q-1$ elements not covered by the linear primes, and no other prime factor.
\end{theorem}
\begin{proof}
$m=1$: there are no primes of degree $<1$, so all $q$ points of $V_1$ are uncovered
(Lemma~\ref{lem:model}); the cost is $q$.

$m=2$: identify $g=c_0+c_1x\in V_2$ with $(c_0,c_1)\in\F_q^2$. The only primes of degree $<2$ are
$x-a$, $a\in\F_q$, and $x-a$ with residue $r$ covers the line $\ell_a(r)=\{c_0+ac_1=r\}$. Lines with
different $a$ have different slopes, so any two meet in exactly one point. If $\ell\ge1$ such lines
are used, each line after the first meets the first one in an already covered point, so at most
$q+(\ell-1)(q-1)$ points are covered and the cost is at least
\[
 \ell+q^2-q-(\ell-1)(q-1)=q^2-1-\ell(q-2)\ \ge\ q^2-1-q(q-2)=2q-1,
\]
using $\ell\le q$ (the count uses only the intersections with the first line, so it does not
depend on the order of the lines); with $\ell=0$ the cost is $q^2\ge2q-1$. For the upper bound take
$f=x^q-x$, $M=\prod_{a\in\F_q}(x-a)\cdot\prod_{c\in\F_q^\times}P_c$, where $P_c$ is an irreducible
factor of $x^q-x+c$. Since $x^q-x+c$ has no root in $\F_q$, $\deg P_c\ge2$, and the $P_c$ are
pairwise distinct because $P_c\mid (c-c')$ would follow otherwise. An element $f+c_0+c_1x$ with
$c_1\ne0$, or with $c_1=0=c_0$, has a root in $\F_q$; the remaining elements are $x^q-x+c$,
$c\ne0$, each divisible by $P_c$. Hence $\omega(M)=2q-1$.

Equality for $q\ge3$: the coefficient $q-2$ is positive, so $\ell=q$, and every line must meet the
union of the previous ones in exactly one point. If three of the lines were not concurrent, the
third would meet the first two in two distinct covered points; hence all $q$ lines pass through one
point $g^*$, i.e.\ $r_{x-a}=g^*(a)$ for all $a$. Since $r_{x-a}=-f(a)$, translating $f$ by $g^*\in V_2$ (which does not
change the coset) gives $(f+g^*)(a)=0$ for all $a$, i.e.\ $x^q-x\mid f+g^*$. Conversely such a
configuration covers $1+q(q-1)$ points and costs $q+(q-1)=2q-1$.
\end{proof}

\begin{remark}
For $q=2$ the characterisation fails: the coefficient $q-2$ vanishes, and extremal configurations exist with $\ell=1$ (one line, two
further primes) as well as with $\ell=2$. Also, linear primes alone
never cover $V_2$: $q$ lines with distinct slopes cover at most $q^2-q+1<q^2$ points.
\end{remark}

\section{Exact small values}
Lemma~\ref{lem:model} turns $K_q(m)$ into a finite set-cover problem. Two independent programs
compute it exactly: an integer program solved with CP-SAT (OR-tools), and a separate exhaustive
branch-and-bound written with its own finite-field and polynomial arithmetic. They agree on:

\begin{center}
\begin{tabular}{lccccc}
\toprule
$q\backslash m$ & 1 & 2 & 3 & 4 & 5\\
\midrule
2 & 2 & 3 & 4 & 6 & 9\\
3 & 3 & 5 & 8 & & \\
4 & 4 & 7 & & & \\
5 & 5 & 9 & & & \\
7 & 7 & 13 & & & \\
\bottomrule
\end{tabular}
\end{center}
($q=4$: branch-and-bound and SageMath only, since the CP-SAT script handles prime $q$.)
A SageMath script verifies the explicit construction of Theorem~\ref{thm:main} for
$q\in\{2,3,4,5,7,8,9\}$ by actual factorisation, and confirms the equality characterisation for
$q\in\{3,4,5\}$ by enumerating all residue vectors. The values $K_2(4)=6$, $K_2(5)=9$, $K_3(3)=8$
exceed $m(q-1)+1$, so the pattern $K_q(m)=m(q-1)+1$ suggested by $m\le2$ does not continue.
No general formula for $m\ge3$ is claimed.

\begin{remark}[related finite geometry]
The bound $K_q(2)\ge2q-1$ is close in spirit to Jamison's theorem and the Brouwer--Schrijver bound on
affine blocking sets \cite{jam} (covering all points of $\mathrm{AG}(2,q)$ but one needs $2q-2$ lines).
Here the lines must have pairwise distinct directions and a missed point costs one prime, which
makes the count elementary.
\end{remark}

\section*{Novelty}
To our knowledge (search of GitHub and arXiv up to 2026-10-09), the polynomial-ring analogue of
Jacobsthal's function has not been studied. Kalmynin and Konyagin \cite{KK} study a
``polynomial analogue'' that is different: shifts of integer polynomial sequences. The function-field
short-interval literature (e.g.\ \cite{powerfree}) treats powerfree polynomials. The community
repository \texttt{jdb19937/jacobsthal} improves the integer bound of \cite{oai} by a power of
$\log k$ and does not consider $\F_q[x]$.

\begin{thebibliography}{9}
\bibitem{oai} OpenAI, \emph{A quadratic bound for Jacobsthal's function}, September 25, 2026,
\url{https://github.com/openai/math} (family 021).
\bibitem{KK} A.~B. Kalmynin and S.~V. Konyagin, \emph{A polynomial analogue of Jacobsthal function},
Izv. Math. \textbf{88} (2024), 225--235; arXiv:2302.00459.
\bibitem{jam} R.~E. Jamison, \emph{Covering finite fields with cosets of subspaces}, J. Combin. Theory
Ser. A \textbf{22} (1977), 253--266; A.~E. Brouwer and A. Schrijver, \emph{The blocking number of an
affine space}, J. Combin. Theory Ser. A \textbf{24} (1978), 251--253.
\bibitem{powerfree} \emph{Short interval results for powerfree polynomials over finite fields},
arXiv:2310.02495.
\end{thebibliography}
\end{document}
